% ============================================================
%  6. Experimental Results
% ============================================================
\section{Experimental Results}
\label{sec:experiments}

\subsection{Datasets}

\paragraph{CoNLL-2003 NER.}
The standard English NER benchmark~\citep{tjong2003conll} contains 14,987
training, 3,466 validation, and 3,684 test sentences drawn from Reuters news
wire, annotated with the PER / ORG / LOC / MISC BIO schema ($|\calY|=9$).
We report entity-level F1 using the \texttt{seqeval} scorer.

\paragraph{Distillation corpus.}
Stages~1 and~2 use WikiText-103-raw-v1~\citep{merity2017pointer}
(${\sim}103$M words of Wikipedia) together with CoNLL-2003 training sentences
as the distillation corpus (${\sim}3.9$M sentences after sentence splitting and
filtering for $\geq 4$ words).

\subsection{Training Setup}

\begin{table}[h]
\centering
\caption{Hyperparameters per stage.}
\label{tab:hparams}
\begin{tabular}{lrrr}
\toprule
\textbf{Hyperparameter} & \textbf{Stage 1} & \textbf{Stage 2} & \textbf{Stage 3} \\
\midrule
Epochs          & 5     & 5     & 10    \\
Batch size      & 32    & 64    & 32    \\
Learning rate   & 5e-5  & 3e-4  & 2e-4  \\
Warmup steps    & 1000  & 500   & 200   \\
Max words       & 64    & 64    & 64    \\
Weight decay    & 0.01  & 0.01  & 0.01  \\
Grad.\ clip     & 1.0   & 1.0   & 1.0   \\
$\beta$ (Eq.~\ref{eq:loss3}) & -- & -- & 0.5 \\
Optimizer       & \multicolumn{3}{c}{AdamW~\citep{loshchilov2019decoupled}} \\
LR schedule     & \multicolumn{3}{c}{Cosine with linear warmup} \\
\bottomrule
\end{tabular}
\end{table}

\paragraph{Noise probabilities (Stage~3).}
$p_{\text{sub}}{=}0.10$, $p_{\text{ins}}{=}0.05$,
$p_{\text{del}}{=}0.05$, $p_{\text{space}}{=}0.02$.

\paragraph{Hardware.}
All models trained on a single NVIDIA Tesla~T4 (16\,GB) via Kaggle
GPU sessions.  Inference benchmarks reported on the same GPU and on
Intel Xeon (single-threaded, CPU).

\subsection{Baselines}

\begin{itemize}
  \item \textbf{BiLSTM-CRF + CharCNN}: Classical character-aware NER
        of~\citet{ma2016end}.
  \item \textbf{DistilBERT}: Knowledge-distilled BERT~\citep{sanh2019distilbert},
        66M params, fine-tuned with a linear NER head.
  \item \textbf{BERT-base}: Standard fine-tuned NER
        baseline~\citep{devlin2019bert}, 110M params.
  \item \textbf{CharBERT}: Character-augmented BERT~\citep{ma2020charbert},
        114M params.
\end{itemize}

\subsection{Noise Evaluation Protocol}

We apply the noise operator $\calN$ (Equations~\ref{eq:nsub}--\ref{eq:nspace})
to the CoNLL-2003 \emph{test} split at three noise intensities:

\begin{table}[h]
\centering
\caption{Noise levels used in evaluation.}
\label{tab:noise_levels}
\begin{tabular}{lcccc}
\toprule
\textbf{Level} & $p_{\text{sub}}$ & $p_{\text{ins}}$ & $p_{\text{del}}$ & $p_{\text{space}}$ \\
\midrule
Low    & 0.05 & 0.02 & 0.02 & 0.01 \\
Medium & 0.10 & 0.05 & 0.05 & 0.02 \\
High   & 0.20 & 0.10 & 0.10 & 0.05 \\
\bottomrule
\end{tabular}
\end{table}

Each evaluation is run five times with different random seeds; we report mean
F1.

\subsection{Main Results}

\begin{table}[h]
\centering
\caption{CoNLL-2003 entity-level F1 on clean test set.}
\label{tab:clean}
\begin{tabular}{lrr}
\toprule
\textbf{Model} & \textbf{Params} & \textbf{F1} \\
\midrule
BiLSTM-CRF + CharCNN~\citep{ma2016end}  & $\sim$8M  & 90.94 \\
DistilBERT~\citep{sanh2019distilbert}   & 66M       & 90.70 \\
BERT-base~\citep{devlin2019bert}        & 110M      & 92.14 \\
CharBERT~\citep{ma2020charbert}         & 114M      & 92.61 \\
\midrule
\textbf{\LightRet{} (ours)}             & \textbf{$\sim$4M} & \textbf{85.7} \\
\bottomrule
\end{tabular}
\end{table}

\begin{table}[h]
\centering
\caption{NER F1 under character noise.  \textbf{Bold} = best per column.}
\label{tab:noisy}
\resizebox{\columnwidth}{!}{%
\begin{tabular}{lrcccc}
\toprule
\textbf{Model} & \textbf{Params} & \textbf{Clean} & \textbf{Low} & \textbf{Medium} & \textbf{High} \\
\midrule
BERT-base   & 110M & 91.25 & 71.89 & 52.06 & 24.13 \\
CharBERT    & 114M & 92.61 & --    & --    & --    \\
DistilBERT  &  66M & 89.93 & 63.19 & 42.96 & 19.40 \\
\midrule
\textbf{\LightRet{}} & \textbf{$\sim$4M}
                     & \textbf{85.7}
                     & \textbf{82.7}
                     & \textbf{78.0}
                     & \textbf{69.2} \\
\bottomrule
\end{tabular}
}
\end{table}

\subsection{Ablation Study}

Table~\ref{tab:ablation} isolates the contribution of each component.

\paragraph{Stage 2 is the most critical component.}
Removing Stage~2 causes the largest single degradation in the entire ablation:
$-14.1$ F1 on clean text and $-12.6$ F1 under medium noise (71.6 vs 85.7).
Without Stage~2, \LightRet{} is initialised randomly and fine-tuned directly
on the CoNLL-2003 NER labels---only ${\sim}15{,}000$ training sentences.
Although \RetVec{} alone provides some signal (71.6 F1, well above random),
the BiGRU--Transformer backbone cannot learn contextual word representations
from NER supervision alone; it requires the large-scale distillation corpus
(WikiText-103 $+$ CoNLL-2003, ${\sim}$millions of sentences) that Stage~2 provides.
The symmetric drop across clean and noisy conditions confirms that Stage~2 builds
the foundational backbone representations, not merely noise robustness.

\paragraph{Stage 1 disproportionately benefits noise robustness.}
Replacing Stage~1 with a direct BERT teacher (subword $\to$ word pooling) yields
a smaller but notable gap: $-$0.9 F1 on clean and $-$1.6 F1 under medium noise
(84.8 vs 85.7 and 76.4 vs 78.0).
The \emph{asymmetry} is informative: Stage~1 improves noisy performance nearly
twice as much as clean performance.
\RetBERT{} processes text at word level---the same granularity as \LightRet{}---so
Stage~2 becomes a clean word-to-word distillation.
Without Stage~1, BERT's subword hidden states must be mean-pooled to word level
before computing the cosine loss, introducing an aggregation mismatch that weakens
the distillation signal and particularly degrades the noise-robust representations
the student would otherwise inherit.
\textbf{Noise augmentation} costs nothing on clean text (85.8 vs 85.7) but is the
sole source of noise robustness---without it, medium-noise F1 collapses by 15.7 points.
Using \textbf{classification loss only} ($\beta{=}1$) degrades both clean ($-1.1$) and
noisy ($-3.8$) F1, confirming that teacher alignment improves generalisation beyond
label supervision alone.
Using \textbf{distillation loss only} ($\beta{=}0$) is catastrophic (2.5 F1)---the
model learns to mimic teacher representations but receives no signal for label assignment.
Finally, replacing \RetVec{} with \textbf{random character embeddings} reduces F1
to 5.2, proving that pretrained character embeddings are foundational, not incidental.

\begin{table}[h]
\centering
\caption{Ablation on CoNLL-2003 (F1).  ``Medium'' = medium noise level.}
\label{tab:ablation}
\resizebox{\columnwidth}{!}{%
\begin{tabular}{lcc}
\toprule
\textbf{Configuration} & \textbf{Clean} & \textbf{Medium} \\
\midrule
\LightRet{} (full, 3-stage)                  & \textbf{85.7} & \textbf{78.0} \\
\quad w/o Stage~1 (skip \RetBERT{})          & 84.8 & 76.4 \\
\quad w/o Stage~2 (direct NER finetune)      & 71.6 & 65.4 \\
\quad w/o noise augmentation in Stage~3      & 85.8 & 62.3 \\
\quad $\calL_3 = \calL_{\text{class}}$ only ($\beta{=}1$) & 84.6 & 74.2 \\
\quad $\calL_3 = \calL_{\text{distill}}$ only ($\beta{=}0$) & 2.5 & 2.9 \\
\quad Replace \RetVec{} with random char emb. & 5.2 & 5.2 \\
\bottomrule
\end{tabular}
}
\end{table}

\subsection{Inference Speed and Memory}

\begin{table}[h]
\centering
\caption{Inference efficiency (single sentence, batch\,=\,1).}
\label{tab:speed}
\begin{tabular}{lrcc}
\toprule
\textbf{Model} & \textbf{Params} & \textbf{GPU (ms)} & \textbf{CPU (ms)} \\
\midrule
BERT-base  & 110M & 2.65 &  7.71 \\
DistilBERT &  66M & 1.51 &  3.50 \\
\textbf{\LightRet{}} & \textbf{$\sim$4M} & \textbf{1.66} & \textbf{10.89} \\
\bottomrule
\end{tabular}
\end{table}
